% ======================================================================
% RAPPORT DE STAGE — TSCAN
% Conversion LaTeX du rapport Markdown (docs/rapport_de_stage.md)
% Compilation : pdflatex rapport_de_stage.tex  (x2 pour la table des matières)
% Encodage : UTF-8 — Moteur recommandé : pdflatex (ou xelatex)
% ======================================================================
\documentclass[11pt,a4paper]{report}

% ---------------------------------------------------------------------
% ENCODAGE ET LANGUE
% ---------------------------------------------------------------------
\usepackage[utf8]{inputenc}
\usepackage[T1]{fontenc}
\usepackage{lmodern}
\usepackage[french]{babel}

% ---------------------------------------------------------------------
% MISE EN PAGE
% ---------------------------------------------------------------------
\usepackage[a4paper,margin=2.5cm]{geometry}
\usepackage{setspace}
\onehalfspacing

% ---------------------------------------------------------------------
% PACKAGES UTILES
% ---------------------------------------------------------------------
\usepackage{graphicx}
\usepackage{booktabs}          % tableaux professionnels
\usepackage{longtable}         % tableaux multi-pages
\usepackage{array}
\usepackage{enumitem}          % contrôle des listes
\usepackage{xcolor}
\usepackage{tcolorbox}         % encadrés
\tcbuselibrary{breakable}
\usepackage{tikz}
\usetikzlibrary{shapes.geometric,arrows.meta,positioning,fit,calc,babel}
\usepackage[hidelinks]{hyperref}
\usepackage{fancyhdr}
\usepackage{titlesec}
\usepackage{caption}
\captionsetup{font=small,labelfont=bf}

% ---------------------------------------------------------------------
% EN-TÊTES ET PIEDS DE PAGE
% ---------------------------------------------------------------------
\pagestyle{fancy}
\fancyhf{}
\fancyhead[L]{\small\leftmark}
\fancyhead[R]{\small\thepage}
\renewcommand{\headrulewidth}{0.4pt}
\fancypagestyle{plain}{%
  \fancyhf{}
  \fancyfoot[C]{\thepage}
  \renewcommand{\headrulewidth}{0pt}
}

% ---------------------------------------------------------------------
% SECTION : ESPACEMENTS ET COULEURS
% ---------------------------------------------------------------------
\definecolor{anticblue}{RGB}{0,60,113}
\definecolor{tscangreen}{RGB}{0,105,92}
\titleformat{\chapter}[display]
  {\normalfont\huge\bfseries\color{anticblue}}
  {\filleft\Large Chapitre \thechapter}{1ex}{\titlerule[1.5pt]\vspace{1ex}\filright}
\titlespacing*{\chapter}{0pt}{20pt}{30pt}
\titleformat{\section}
  {\normalfont\Large\bfseries\color{anticblue}}{\thesection}{0.8em}{}
\titleformat{\subsection}
  {\normalfont\large\bfseries}{\thesubsection}{0.8em}{}

% ---------------------------------------------------------------------
% COMMANDES UTILES
% ---------------------------------------------------------------------
\newcommand{\tscan}{\textsc{Tscan}}
\newcommand{\todo}[1]{\textcolor{red}{\textbf{[#1]}}}
\newcommand{\code}[1]{\texttt{#1}}

% Boîte "À COMPLÉTER" bien visible
\newtcolorbox{acompleter}[1][]{%
  colback=red!4!white, colframe=red!60!black, breakable,
  title={\textbf{À compléter}}, #1
}

% Boîte "À retenir"
\newtcolorbox{aretenir}[1][]{%
  colback=tscangreen!5!white, colframe=tscangreen, breakable,
  title={\textbf{Point clé}}, #1
}

% =====================================================================
\begin{document}

% =====================================================================
% PAGE DE GARDE
% =====================================================================
\begin{titlepage}
\begin{center}

{\Large \textbf{RÉPUBLIQUE DU CAMEROUN}}\\
{\small Paix -- Travail -- Patrie}\\[0.3cm]
{\normalsize \textbf{AGENCE NATIONALE DES TECHNOLOGIES DE L'INFORMATION\\ ET DE LA COMMUNICATION (ANTIC)}}\\[1.2cm]

\rule{\textwidth}{1.5pt}\\[0.4cm]
{\Huge \bfseries Rapport de stage}\\[0.6cm]
{\LARGE \bfseries \textsc{Tscan}}\\[0.5cm]
{\large Plateforme d'analyse, de corrélation et de validation de vulnérabilités}\\[0.4cm]
\rule{\textwidth}{1.5pt}\\[1.5cm]

\begin{tabular}{rl}
  \textbf{Stagiaire :} & KEN DANIEL Takou \\[0.2cm]
  \textbf{Formation :} & Cycle de formation en génie civil \\[0.2cm]
  \textbf{Organisme d'accueil :} & ANTIC --- Agence Nationale des Technologies \\ 
   & de l'Information et de la Communication \\[0.2cm]
  \textbf{Période :} & \todo{20/07/2026 $\rightarrow$ 18/09/2026} \\[0.2cm]
  \textbf{Encadrant :} & \todo{Nom de l'encadrant} \\[0.2cm]
  \textbf{Service d'accueil :} & \todo{Département / équipe} \\
\end{tabular}\\[1.5cm]

\begin{tcolorbox}[colback=anticblue!4!white, colframe=anticblue, width=0.85\textwidth]
\centering \textit{« Réduire le temps de tri des alertes de sécurité et améliorer la fiabilité des constats, en croisant les résultats des scanners externes avec un moteur de validation interne. »}
\end{tcolorbox}

\vfill
{\large Version 1.0 --- \today}
\end{center}
\end{titlepage}

% =====================================================================
% PAGE DE REMERCIEMENTS (À PERSONNALISER)
% =====================================================================
\chapter*{Remerciements}
\addcontentsline{toc}{chapter}{Remerciements}

Je tiens à remercier \todo{Nom de l'encadrant}, mon encadrant de stage, pour sa
disponibilité, ses orientations et la confiance accordée tout au long de ce projet. Mes
remerciements vont également à l'ensemble de l'équipe du \todo{département d'accueil} de
l'ANTIC pour la qualité de l'accueil et les conditions de travail offertes.

Enfin, je remercie la direction de ma formation en génie civil pour avoir autorisé et
encadré ce stage, ainsi que toutes les personnes qui, de près ou de loin, ont contribué à
la réussite de ce travail.

% =====================================================================
% TABLE DES MATIÈRES
% =====================================================================
\tableofcontents

% =====================================================================
% SYNTHÈSE (RÉSUMÉ)
% =====================================================================
\chapter*{Synthèse}
\addcontentsline{toc}{chapter}{Synthèse}

Ce rapport présente le travail réalisé au cours de mon stage à l'ANTIC, consacré au
développement de \tscan{} : une plateforme d'analyse, de corrélation et de validation de
vulnérabilités de sécurité pour applications web. Le stage couvre l'intégralité du cycle
produit --- étude de l'existant, cahier des charges, architecture, développement, tests,
validation expérimentale --- sur la base d'un cahier des charges de 16 chapitres et d'un
planning en 11 semaines.

La particularité du projet tient à son \textbf{double enjeu} : construire un outil
opérationnel (import multi-formats, scan actif non destructif, corrélation, reporting), et
répondre à un problème concret des équipes de sécurité : \textbf{les scanners produisent
beaucoup de faux positifs}, et la fiabilisation de leurs résultats exige des méthodes
structurées. \tscan{} attaque ce problème par trois leviers : le \textbf{pré-filtrage
contextuel}, la \textbf{corrélation multi-sources} et la \textbf{validation active non
destructive} (ré-observation).

Le présent document décrit les choix techniques, les livrables, la validation
expérimentale mesurée sur cibles réelles et autorisées, ainsi que les limites assumées du
produit au terme du stage.

% =====================================================================
% LISTE DES FIGURES ET TABLEAUX (optionnel mais recommandé)
% =====================================================================
\cleardoublepage
\addcontentsline{toc}{chapter}{Liste des figures et des tableaux}
\listoffigures
\listoftables

% =====================================================================
\chapter{Introduction}

Ce rapport présente le travail réalisé au cours de mon stage à l'ANTIC, consacré au
développement de \tscan{} : une plateforme d'analyse, de corrélation et de validation de
vulnérabilités de sécurité pour applications web. Le stage couvre l'intégralité du cycle
produit --- étude de l'existant, cahier des charges, architecture, développement, tests,
validation expérimentale --- sur la base d'un cahier des charges de 16 chapitres et d'un
planning en 11 semaines.

La particularité du projet tient à son \textbf{double enjeu} : construire un outil
opérationnel (import multi-formats, scan actif non destructif, corrélation, reporting), et
répondre à un problème concret des équipes de sécurité : \textbf{les scanners produisent
beaucoup de faux positifs}, et la fiabilisation de leurs résultats exige des méthodes
structurées. \tscan{} attaque ce problème par trois leviers : le \textbf{pré-filtrage
contextuel}, la \textbf{corrélation multi-sources} et la \textbf{validation active non
destructive} (ré-observation).

Le présent document décrit les choix techniques, les livrables, la validation
expérimentale mesurée sur cibles réelles et autorisées, ainsi que les limites assumées du
produit au terme du stage.

% =====================================================================
\chapter{Organisation d'accueil}

\section{L'ANTIC}

L'\href{https://www.antic.cm}{ANTIC} (Agence Nationale des Technologies de l'Information
et de la Communication) est l'agence gouvernementale camerounaise chargée de la promotion
et de la sécurisation des technologies de l'information. Ses missions couvrent notamment
l'administration des ressources Internet du pays (domaine national «~.cm~», infrastructure
DNS et registre), le développement des usages numériques et la cyberdéfense nationale. Elle
héberge le \textbf{CERT/CCTAL}, l'équipe de réponse aux incidents de sécurité informatique
du Cameroun.

\section{Mon rôle et mon encadrement}

J'ai été accueilli au sein du \todo{département/équipe} et encadré par
\todo{Nom de l'encadrant}. Mon travail a suivi la logique MVP (Minimum Viable Product)
définie au chapitre 5 du cahier des charges : livrer d'abord le cœur du produit, puis
étendre par blocs hebdomadaires testés et vérifiés. Les points d'avancement hebdomadaires
ont été consignés dans une checklist projet (\code{docs/checklist\_projet.md}) comparant
avancement réel et planning, et le README reflète à tout instant l'état réel du dépôt.

% =====================================================================
\chapter{Contexte et problématique}

\section{Le contexte}

Les organisations sont soumises à une pression croissante : surface d'attaque web étendue,
composants tiers omniprésents (CMS, bibliothèques JavaScript), exigences de conformité.
Les scanners de vulnérabilités (Nuclei, OWASP ZAP, Nessus, OpenVAS) sont devenus
indispensables, mais leur sortie brute présente trois faiblesses :

\begin{itemize}[leftmargin=1.5em]
  \item \textbf{formats hétérogènes} : chaque outil a son modèle de rapport, rendant
        difficile une vue consolidée ;
  \item \textbf{faux positifs} : une alerte détectée n'est pas toujours une vulnérabilité
        réelle --- le filtrage manuel des alertes est le premier poste de charge des
        analystes ;
  \item \textbf{confiance non explicable} : un score 5/10 sans justification n'aide pas à
        prioriser.
\end{itemize}

\section{La problématique}

\begin{tcolorbox}[colback=anticblue!4!white, colframe=anticblue,
                  title={\textbf{Problématique du stage}}]
Comment réduire le temps de tri des alertes et améliorer la fiabilité des constats, en
cross-checkant les résultats des scanners externes avec un moteur de validation interne ?
\end{tcolorbox}

\tscan{} répond par un modèle pivot unique (\code{Finding}), un pré-filtrage contextuel,
une corrélation multi-sources avec scoring de confiance \textbf{explicable}, et une
\textbf{ré-observation active} : chaque constat importé peut être re-validé par une seconde
observation non destructive exécutée par le moteur interne.

% =====================================================================
\chapter{Objectifs et périmètre du stage}

\section{Objectif général}

Concevoir et développer \tscan{}, plateforme d'analyse, de corrélation et de validation de
vulnérabilités web, conformément au cahier des charges (16 chapitres) et à son MVP
(chapitre 5), avec un niveau de qualité industrialisable : tests automatisés, style
vérifié, traçabilité des décisions (checklist + README reflétant l'état réel).

\section{Objectifs spécifiques (par bloc hebdomadaire)}

\begin{table}[h!]
\centering
\small
\begin{tabular}{@{}cp{9.8cm}p{3.2cm}@{}}
\toprule
\textbf{Bloc} & \textbf{Objectif} & \textbf{État final} \\
\midrule
S1--S2 & Étude de l'existant, cahier des charges, architecture prévisionnelle & Terminé \\
S3 & Squelette technique (venv, pyproject, ruff, pytest, modèles, SQLite, CLI Typer) & Terminé \\
S4 & Import Nuclei/ZAP vers le modèle pivot \code{Finding} & Terminé \\
S5--S6 & Corrélation, scoring explicable, statuts, historique --- scénario A en CLI & Terminé \\
+ & Bloc connaissances CVE/NVD + CISA KEV + références CWE + alias CPE & Terminé \\
S7 & Reconnaissance (statut, TLS, technologies), périmètre, journalisation + détections passives & Terminé \\
S8 & Détections actives non destructives (XSS, CSRF, SQLi error-based\ldots) + import Nessus/OpenVAS & Terminé \\
S9 & Confirmation active RF-23 + reporting (HTML/Markdown) + scénario B & Terminé \\
S10 & Application desktop PySide6 (liste filtrable, détail, correction, import/scan/rapport) & Terminé \\
S11 & Durcissement, empaquetage, démonstration & Terminé \\
\bottomrule
\end{tabular}
\caption{Objectifs hebdomadaires et état final.}
\end{table}

\section{Périmètre du MVP}

Le MVP couvre l'import des formats \textbf{Nuclei (JSONL), OWASP ZAP (JSON),
Nessus/OpenVAS (XML)}, la corrélation multi-sources sur cible commune, le scoring de
confiance explicable, les statuts avec correction manuelle et historique, un \textbf{moteur
de scan actif autorisé et non destructif}, la ré-observation (corroboration), le reporting
HTML/Markdown, une CLI complète et une application desktop. \textbf{Hors périmètre} : PDF,
SSRF/désérialisation (familles à risque, ES-04), gestion multi-utilisateurs,
authentification applicative.

% =====================================================================
\chapter{Étude de l'existant}

Avant le développement, une étude de l'existant a été conduite sur quatre axes :

\begin{itemize}[leftmargin=1.5em]
  \item \textbf{Scanners de vulnérabilités} : Nuclei (templating communautaire), OWASP ZAP
        (proxy d'analyse avec scanner actif), Nessus/OpenVAS (scanners infra). Aucun ne
        fait nativement de corrélation inter-outils avec fiabilisation.
  \item \textbf{Plateformes d'agrégation} (modèle «~vuln management~») : se limitent
        souvent à la consolidation, sans validation active.
  \item \textbf{Sources de connaissances} : NVD/CVE, CISA KEV (catalogue des exploitations
        avérées), CWE pour la classification des faiblesses.
  \item \textbf{Formats de rapport} : JSONL (Nuclei), JSON (ZAP), XML \code{.nessus}
        (Nessus/OpenVAS).
\end{itemize}

\begin{aretenir}
Il existe des outils excellents à chaque étage, mais \textbf{aucun chaînon mature qui
corrèle plusieurs sources sur une même cible, score la confiance de façon explicable, puis
re-valide activement les constats importés}. C'est la niche que \tscan{} occupe.
\end{aretenir}

% =====================================================================
\chapter{Choix méthodologiques}

\section{Démarche MVP incrémentale}

Le développement a suivi une logique MVP : chaque bloc hebdomadaire livre des
fonctionnalités testées et vérifiées, consignées dans une checklist projet avec comparaison
avancement réel / planning. Le README reflète l'état réel du dépôt à chaque instant.

\section{Assurance qualité}

\begin{itemize}[leftmargin=1.5em]
  \item \textbf{343 tests automatisés} (pytest), tous verts ;
  \item \textbf{ruff} (linting) propre sur \code{src/}, \code{tests/} et \code{rules/} ;
  \item \textbf{1 test automatisé minimum par module critique} ;
  \item audit des dépendances (\code{pip-audit} via \code{scripts/audit\_deps.py}, ES-12) ;
  \item migration de données idempotente à l'initialisation de la base (RF-12).
\end{itemize}

\section{Pile technologique}

\begin{table}[h!]
\centering
\small
\begin{tabular}{@{}llp{7.5cm}@{}}
\toprule
\textbf{Couche} & \textbf{Choix} & \textbf{Justification} \\
\midrule
Langage & Python $\geq$ 3.11 & Écosystème sécurité, httpx, rapidité de développement \\
Bibliothèque cœur & modules \code{tscan\_core} & Partagée CLI + GUI, testable sans écran \\
CLI & Typer & Complétion, sous-commandes, ergonomie \\
Desktop & PySide6 & Qt pour Python, LGPL, opérations longues en QThread \\
Persistance & SQLite + SQLAlchemy & Base fichier unique, migrations additives, zéro infra \\
HTTP & httpx & Client borné (timeout, redirections limitées), brotli/gzip \\
Règles & YAML versionnés & Aucune règle codée en dur, chargeur dynamique \\
Tests & pytest + serveur labo local & Intégration sans réseau externe \\
Empaquetage & PyInstaller + Inno Setup & Exécutable Windows (S11) \\
\bottomrule
\end{tabular}
\caption{Pile technologique retenue.}
\end{table}

\section{Structure du dépôt}

\begin{verbatim}
src/
  tscan_core/       Bibliothèque cœur (import, corrélation, détection, reporting, scan)
  tscan_cli/        Interface en ligne de commande (Typer)
  tscan_gui/        Application desktop (PySide6)
rules/              54 règles de détection versionnées (YAML)
knowledge/          Technologies, alias CPE, référence CWE, signatures de faux positifs
tests/              343 tests pytest + labo local (lab_server.py) + fixtures réelles
docs/               Cahier des charges (PDF), checklist projet, rapport de stage
\end{verbatim}

% =====================================================================
\chapter{Architecture technique}

\section{Vue d'ensemble}

\begin{figure}[h!]
\centering
\begin{tikzpicture}[
  node distance=0.55cm and 0.4cm,
  layer/.style={rectangle, rounded corners=3pt, draw=anticblue, thick,
                fill=anticblue!6, align=center, inner sep=6pt},
  core/.style={rectangle, rounded corners=3pt, draw=tscangreen, thick,
               fill=tscangreen!6, align=center, inner sep=6pt},
  data/.style={rectangle, rounded corners=3pt, draw=black!60, thick,
               fill=black!4, align=center, inner sep=6pt},
  arr/.style={-{Stealth[length=3mm]}, thick, black!70}
]

% Couche présentation
\node[layer, text width=11.5cm] (gui) {\textbf{tscan\_gui (PySide6)}\\
  MainWindow · FilterBar · FindingsTable · FindingDetailPanel\\
  Dialogs · Workers (QThread)};
\node[layer, below=of gui, text width=11.5cm] (cli) {\textbf{tscan\_cli (Typer)}\\
  Commandes : import · scan · correlate · report · correct · lookup-cve · update-kev};

% Coeur
\node[core, below=of cli, text width=11.5cm] (core) {\textbf{tscan\_core (bibliothèque cœur)}\\
  modèle pivot Finding $\rightarrow$ pré-filtrage $\rightarrow$ corrélation $\rightarrow$
  scoring $\rightarrow$ statuts/historique $\rightarrow$ reporting HTML/MD\\[2pt]
  \begin{tabular}{cc}
    \begin{tabular}[t]{@{}l@{}}\textbf{Scan actif}\\ recon $\rightarrow$ règles\\
      $\rightarrow$ ré-observation\end{tabular} &
    \begin{tabular}[t]{@{}l@{}}\textbf{Base de connaissances}\\ NVD/CVE · KEV · CWE/CPE\\
      (cache local, à la demande)\end{tabular}
  \end{tabular}};

% Importeurs
\node[data, left=0.5cm of core.north west, anchor=north west, yshift=-1.2cm,
      text width=2.6cm] (imports)
  {\textbf{Importeurs}\\ Nuclei (JSONL)\\ ZAP (JSON)\\ Nessus/OpenVAS (XML)};

% Stockage
\node[data, below=of core, text width=11.5cm] (db)
  {\textbf{SQLite} --- Scans · Findings · Evidence · StatusHistory · RequestHistory};

% Flèches
\draw[arr] (gui) -- (cli);
\draw[arr] (cli) -- (core);
\draw[arr] (imports) -- (core);
\draw[arr] (core) -- (db);

\end{tikzpicture}
\caption{Architecture en couches de \tscan{}.}
\end{figure}

\section{Le modèle pivot \code{Finding}}

Tout résultat --- importé ou produit par le scan actif --- converge vers un modèle unique :
cible, emplacement (\code{matched\_at}), catégorie, titre, gravité, statut, score de
confiance, preuves (\code{Evidence}), résultat brut conservé intact (ES-07), historique de
statut (\code{StatusHistory}, ES-06). C'est ce pivot qui rend la corrélation et le scoring
possibles entre sources hétérogènes.

\section{Le moteur de règles (RF-13)}

\textbf{Aucune règle codée en dur} : 54 règles YAML dans \code{rules/}, une par famille de
détection, chargées dynamiquement. Chaque règle déclare identifiant, catégorie, gravité,
conditions, méthode de validation, preuves attendues, recommandations sourcées
(OWASP/CWE) et version. Le moteur de checks déclaratifs (section \code{checks:}) exécute
les vérifications passives (\code{header\_absent}, \code{clickjacking},
\code{path\_status}\ldots) ; les familles actives (XSS, SQLi\ldots) implémentent leurs
sondes en Python, bornées par le module \code{probe.py}.

\section{Le scan actif (RF-15--24, ES-01--05)}

Le pipeline d'un scan actif autorisé :

\begin{enumerate}[leftmargin=1.8em]
  \item \textbf{Reconnaissance} (RF-15) : requête racine via client HTTP borné (timeout, 3
        redirections max, pacing, rotation d'User-Agent), sonde TLS, fingerprinting
        (RF-16) des technologies (PHP, WordPress, jQuery\ldots) via
        \code{knowledge/technologies.yaml} ;
  \item \textbf{Crawl} borné (profondeur et pages plafonnées, pages HTML uniquement pour
        les sondes) ;
  \item \textbf{Détections} selon le périmètre (\code{--tests}, RF-24) : 27 familles
        couvrant les en-têtes de sécurité, clickjacking, BAC, composants vulnérables
        (cache CVE local, ES-09), XSS réfléchi, CSRF, SQLi error-based, SSRF \emph{non
        exécuté}, SSTI, fichiers sensibles, listing, CORS, TLS faible, cookies, WAF, CSP,
        SRI, inclusions cross-domain, timestamps, en-têtes passifs parité ZAP\ldots ---
        toutes \textbf{GET-only bénignes} (ES-03) ;
  \item \textbf{Ré-observation} (RF-23) : re-vérification active de chaque constat
        \code{Probable} ; un fait reproduit \textbf{renforce la confiance (plafonné
        0,85)} sans jamais changer le statut ; un fait contredit place le constat en
        \code{Potentiel faux positif} ; une sonde en échec laisse le constat tel quel
        (ES-05). \textbf{Le moteur ne confirme jamais} : «~Confirmée~» (0,95) est un
        verdict d'analyste (RF-12).
\end{enumerate}

\section{La corroboration (ré-observation des imports)}

C'est la fonctionnalité au cœur de la problématique : réconcilier un \textbf{rapport
importé} (par ex. ZAP) avec une \textbf{observation fraîche} du moteur interne :

\begin{itemize}[leftmargin=1.5em]
  \item \textbf{passe 1 --- URL exacte} : catégorie + emplacement normalisé identiques ;
  \item \textbf{passe 2 --- titre + catégorie + hôte} : même vulnérabilité, instance
        différente (l'URL précise diffère, la faiblesse est la même sur le même hôte) ;
  \item \textbf{table de synonymes FR$\leftrightarrow$EN} : les rapports ZAP en français
        («~Incompatibilité de charset~») rejoignent les titres du moteur («~Charset
        Mismatch~») ;
  \item \textbf{catégorie «~other~» en joker} : un titre sans indice fiable est rapproché
        par titre seul.
\end{itemize}

Le taux de corroboration mesuré sur le cas d'étude est présenté au chapitre 9.

\section{La base de connaissances}

Accès \textbf{à la demande, avec cache local} (aucun réseau pendant un scan, ES-09) :
NVD/CVE (requête CVE par CVE), CISA KEV (catalogue des vulnérabilités exploitées,
synchronisation complète via \code{tscan update-kev}), référence CWE minimale et table
d'alias CPE vérifiée contre l'API NVD. Garde-fous d'intégrité (ES-11) : taille maximale des
réponses + validation de format CVE avant intégration.

% =====================================================================
\chapter{Réalisation : du squelette au moteur de scan actif}

\section{Semaines 3--4 : socle et modèle pivot}

Mise en place du squelette (venv, \code{pyproject.toml}, ruff, pytest, modèles SQLAlchemy
\code{Finding}/\code{Scan}/\code{Rule}/\code{Evidence}, base SQLite à migrations additives,
CLI Typer). Puis le bloc import : parseurs \textbf{Nuclei (JSONL)} et \textbf{OWASP ZAP
(JSON)} vers le modèle pivot, avec conservation du résultat brut (ES-07) et tests sur jeux
de données réels.

\section{Semaines 5--6 : corrélation, scoring, statuts}

Pré-filtrage contextuel (RF-08), corrélation multi-sources (RF-09), moteur de scoring
\textbf{explicable} (RF-10) --- chaque score dérive de la règle appliquée, du bonus
multi-sources et de la re-vérification ---, consultation des preuves (RF-11), correction
manuelle avec historique (RF-12). \textbf{Jalon : scénario A démontrable en CLI}.

\section{Bloc connaissances (entre S6 et S7)}

CVE/NVD à la demande avec cache local, CISA KEV, référence CWE, alias CPE --- commandes
\code{tscan lookup-cve} et \code{tscan update-kev}.

\section{Semaine 7 : reconnaissance et détections passives}

Reconnaissance (client HTTP borné, TLS, fingerprinting), périmètre (cible autorisée,
profondeur, durée, types de test --- RF-24/ES-01/02), journalisation systématique (ES-05),
mode sécurisé par défaut (ES-03). Détections passives : en-têtes de sécurité manquants,
clickjacking, composants vulnérables (correspondance version $\rightarrow$ CVE sur cache
local), broken access control basique. Vérification sur \textbf{labo local} puis sur cible
réelle autorisée.

\section{Semaine 8 : détections actives et validation active}

XSS réfléchi (charge bénigne à nonce), formulaires POST sans jeton CSRF, SQLi error-based
(charges \code{' --}/\code{' \#} + marqueurs d'erreur), plus quatre familles
complémentaires : fichiers sensibles, listing de répertoire, CORS permissif, TLS faible.
Import Nessus/OpenVAS. \textbf{Confirmation active RF-23} puis \textbf{jalon de sécurité :
scénario B} (reconnaissance $\rightarrow$ familles $\rightarrow$ validation $\rightarrow$
scoring $\rightarrow$ rapport).

\section{Semaine 9 : reporting et finalisation CLI}

Générateur de rapport (résumé exécutif par gravité/statut + détails techniques, preuves,
recommandations liées à des sources vérifiables OWASP/CWE --- RF-27/28/29), exports
\textbf{HTML et Markdown}, commande \code{tscan report}, tests d'intégration bout en bout,
revue d'architecture et revue de code. \textbf{Défaut corrigé} au passage : la
désambiguïsation des règles partageant une catégorie (\code{scoring.\_resolve\_rule}).

\section{Semaine 10 : application desktop}

Fenêtre principale (liste filtrable statut/gravité/cible/recherche --- RF-11), vue détail
(preuves, score, historique), correction manuelle avec raison consignée (ES-06),
\code{ImportDialog}, \code{ScanDialog} (périmètre + autorisation), \code{ReportDialog}
(HTML/Markdown). Opérations longues en \textbf{QThread} (RNF-03) : l'interface reste
réactive. Logique de présentation testée sans écran (viewmodel). Un défaut de refactor
(fichiers invalides) a été corrigé le 27/08/2026, la fenêtre a été ouverte et vérifiée en
headless.

\section{Semaine 11 : durcissement et correctifs}

\begin{itemize}[leftmargin=1.5em]
  \item \textbf{Corrélation cohérente} : \code{run\_correlation} n'inclut plus les
        résultats du scan actif ni les constats déjà Confirmés --- corréler ne rétrograde
        plus une confirmation humaine ;
  \item \textbf{Progression temps réel} : événements de progression émis par le moteur
        (\code{tscan\_core.scan.progress}), panneau dans la GUI (barre + journal, comme
        ZAP) ;
  \item \textbf{Migration des données} : les anciens \code{Confirmée} auto-posés par le
        moteur sont rétrogradés en \code{Probable} (idempotent, historique conservé) ;
  \item \textbf{Faux positifs démontrés} : test d'intégration avec serveur qui expose puis
        protège \code{/admin/} $\rightarrow$ constat en \code{Potentiel faux positif}.
\end{itemize}

\section{Travaux de fin de stage (septembre 2026)}

Cette dernière période a concentré l'effort sur la \textbf{fiabilisation et la performance
du moteur}, avec un cycle complet diagnostic $\rightarrow$ correctif $\rightarrow$
validation.

\subsection{Parité passive avec OWASP ZAP}

Analyse croisée des alertes ZAP réelles et des constats du moteur : ajout d'un bundle de
\textbf{16 règles passives} (\code{zap\_passive.py}) aux titres exacts ZAP (Content-Type
Missing, Suspicious Comments, Modern Web Application, Cache-control, Retrieved from Cache,
Permissions-Policy, X-Frame-Options, Server Header Leak, X-AspNet-Version,
X-ChromeLogger-Data, X-Debug-Token, Reverse Tabnabbing, Weak Authentication,
Hash/Base64/Source Code Disclosure) et d'une règle \textbf{XSS attribut} (RF : «~User
Controllable HTML Element Attribute (Potential XSS)~», ZAP 10031). Câblage complet :
famille \code{zap-passives}, 17 règles YAML, comptages de tests mis à jour.

\subsection{Normalisation des catégories FR/EN}

Extension de \code{normalize\_category} (mots-clés timestamp, content-type, modern web app,
cache, server header, weak auth, base64/hash/source code disclosure, tabnabbing,
x-aspnet/chromelogger/debug-token, auth request, session management) ; correctif du
matching par mot entier pour les sigles courts (\code{rce} ne doit plus matcher
«~r\textbf{esou}rce~») ; \textbf{table de synonymes FR$\leftrightarrow$EN} pour la passe
titre de la corroboration. Correctif de régression : restauration du mapping Nessus
(\code{cve-}, \code{.env}) via seuil de sous-chaîne $\geq$ 4.

\subsection{Décompression Brotli --- cause racine d'un échec silencieux}

Le client annonçait \code{Accept-Encoding: gzip, deflate, br} sans le paquet
\code{brotli} : sur les serveurs compressant en Brotli (LiteSpeed), \textbf{tout le HTML
analysé était un blob illisible} --- 0 détection passive, crawl bloqué à 1 page. Correctif :
\code{brotli} installé et ajouté aux dépendances. Vérifié en conditions réelles : 39
scripts, 4 formulaires et 250 liens soudain visibles sur une page WordPress ; détecteurs
passant de 5 à 10 constats sur la seule page racine, incluant exactement les alertes ZAP
manquantes à la corroboration.

\subsection{Performance du scan actif (durée prévisible)}

Le déblocage du crawl sur un CMS a fait exploser le volume de sondes : 17 pages
$\times$ 58 paramètres $\times$ 2 charges $\approx$ \textbf{2\,000 requêtes} rien que pour
la détection SQLi. Correctifs : fonction \code{probe\_fetch} (1 tentative, timeout 5~s)
adoptée par les familles volumineuses ; plafonnement du nombre de paramètres sondés
(\code{MAX\_PARAMS\_PER\_SCAN = 10}, paramètre de recherche toujours couvert) ; cadence
réduite de la re-vérification ; arrêt anticipé des sondes big-redirect après échecs
consécutifs. Volume SQLi ramené à $\sim$240 requêtes ; objectif de durée : $\sim$5--15 min.

\subsection{Message d'erreur explicite en cas de blocage anti-bot}

Lorsqu'une cible dropperait les connexions en silence (WAF/anti-bot), le message d'erreur
distingue désormais ce cas d'une panne réseau et invite à réessayer plus tard --- retour
utilisateur honnête, comme le prévoit l'ES-05.

\subsection{Démarche anti-faux-positifs (sentinelle, signatures, cohérence)}

Le faux positif «~panneau d'administration accessible sans authentification~» observé sur
la cible réelle a servi de fil conducteur à un lot de fiabilisation du check
\code{path\_status} (RF-20) :

\begin{enumerate}[leftmargin=1.8em]
  \item \textit{Redirections vers une page de connexion.} La sonde suivait la redirection
        \code{302 /wp-admin/ $\rightarrow$ /wp-login.php} et ne regardait que le statut
        final (200) : le 200 de la page de connexion était lu comme celui du panneau. Le
        verdict distingue désormais la destination (\code{Location}) et le trajet (URL
        finale $\neq$ URL demandée).
  \item \textit{Sentinelle anti «~soft-404~» (contrôle négatif).} Certains serveurs
        répondent 200 avec une page générique pour n'importe quel chemin. Avant les sondes
        de chemins, une requête vers un chemin inexistant aléatoire
        (\code{/tscan-calib-<jeton>}) établit le «~bruit de fond~» du serveur ; une
        réponse de sonde indiscernable du contrôle (réflexion du jeton, corps quasi
        identique) ne produit pas de constat --- un 200 ne prouve plus rien. Le contrôle
        est partagé entre détection (orchestrateur) et re-vérification (RF-23), où il
        ajoute une note «~fait décisif non démontrable~» plutôt qu'une dégradation
        silencieuse.
  \item \textit{Signatures de faux positifs en connaissances.} Les marqueurs de pages
        d'authentification (WordPress, Drupal, Django, pages francophones
        \code{connexion}\ldots) ont quitté le code pour
        \code{knowledge/fp\_signatures.yaml}, au même principe que les règles YAML et les
        autres tables du dossier \code{knowledge/} : la table s'étend sans toucher au
        code, avec des consignes anti-collisions documentées («~auth~» attraperait
        \code{/author/john-doe}) et un test qui garantit qu'aucun marqueur ne capture une
        ressource réellement sondée.
\end{enumerate}

Le compromis retenu, fidèle au RF-12 : en cas de doute, le moteur \textbf{s'abstient et
documente} (note pour revue analytique) au lieu de conclure --- il ne classe jamais un
constat en faux positif sur une simple incertitude de mesure.

\subsection{Publication du projet}

Le code a été publié sur GitHub (\code{github.com/Takou237/tscan2}), avec historique de
commits, \code{.gitignore} protégeant les données de scan locales (bases \code{*.db}
jamais versionnées) et exclusion des scripts d'analyse personnels contenant des chemins
locaux.

% =====================================================================
\chapter{Validation expérimentale et mesure de corroboration}

\section{Protocole}

Pour mesurer la capacité du moteur à \textbf{reproduire les constats d'un scanner tiers},
le protocole suivant a été appliqué à un site réel autorisé (boutique
WordPress/WooCommerce) :

\begin{enumerate}[leftmargin=1.8em]
  \item \textbf{Import} d'un rapport ZAP réel (26 alertes) vers le modèle pivot ;
  \item \textbf{Scan actif} \tscan{} sur la même cible (mêmes contraintes de bonté :
        GET-only) ;
  \item \textbf{Corroboration} : rapprochement import $\leftrightarrow$ moteur (passe URL
        exacte puis passe titre+catégorie+hôte avec synonymes FR/EN) ;
  \item \textbf{Mesure} du taux : corroborés / alertes importées.
\end{enumerate}

La cible a été scannée avec autorisation, à faible volume, en mode non destructif
(ES-01/02/03). Les constats importés non reproduits sont classés \code{non reproductible}
(revue analytique manuelle) ou \code{faux positif potentiel} lorsque la re-vérification
contredit explicitement le fait.

\section{Résultats}

\begin{table}[h!]
\centering
\begin{tabular}{@{}lc@{}}
\toprule
\textbf{Étape} & \textbf{Taux de corroboration} \\
\midrule
Avant le lot de détecteurs (situation initiale, 16/26) & 61,5~\% \\
Après parité ZAP + normalisation FR/EN (17/26 hors ligne) & 65,4~\% \\
Après correctif Brotli + scan actif complet (20/26) & \textbf{76,9~\%} \\
\bottomrule
\end{tabular}
\caption{Taux de corroboration mesuré sur le cas d'étude.}
\end{table}

\textbf{Analyse des non-corroborés restants (6/26)} : pour la moitié, l'alerte ZAP portait
sur des ressources que le moteur \tscan{} ne crawle volontairement pas (fichiers images
\code{.png}/\code{.jpg} ciblés par ZAP pour l'absence de CSP) ; pour les autres, le site a
commencé à \textbf{filtrer les connexions du scanner} en cours de scan (WAF anti-bot) ---
comportement documenté au chapitre 10. Ces limites sont méthodologiquement significatives :
\textbf{les deux outils sont sujets au même blocage}, ce qui relativise l'écart constaté.

\textbf{Indicateurs qualité} : 343 tests verts, ruff propre, un test minimum par module
critique, migration de données vérifiée sur la base réelle. L'effet des correctifs
anti-faux-positifs de la semaine 12 est mesuré en 9.4.

\section{Comparaison ZAP vs Tscan sur le cas d'étude}

\begin{table}[h!]
\centering
\small
\begin{tabular}{@{}p{10.5cm}c@{}}
\toprule
\textbf{Alerte ZAP importée} & \textbf{Tscan reproduit ?} \\
\midrule
En-têtes manquants (HSTS, XCTO, XFO, Permissions-Policy, CSP) & oui \\
Cookies sans HttpOnly/SameSite, Cookie Poisoning & oui \\
X-Powered-By divulgué, Timestamps Unix, Suspicious Comments, Modern Web App & oui \\
Cache-control / Retrieved from Cache & oui \\
Authentication / Session Management Identified & oui (règles passives parité ZAP) \\
Sub Resource Integrity Missing, Cross-Domain JS & oui (après correctif Brotli) \\
Absence de jetons Anti-CSRF & oui (formulaires POST analysés) \\
CSP Not Set sur image PNG, charset, XSS attribut, Sensitive Info in URL & partiel / non \\
\bottomrule
\end{tabular}
\caption{Correspondance alertes ZAP / détections \tscan{} sur le cas d'étude.}
\end{table}

\section{Effet des correctifs anti-faux-positifs (semaine 12)}

Le lot de fiabilisation décrit en 8.9 (verdict sur la destination des redirections,
sentinelle anti soft-404, signatures déclaratives) a été mesuré sur la cible d'étude en
conditions réelles (17/09/2026, scan actif \#5, autorisation inchangée, même adresse IP) :

\begin{table}[h!]
\centering
\small
\begin{tabular}{@{}p{5.6cm}p{4.6cm}p{4.6cm}@{}}
\toprule
\textbf{Observation} & \textbf{Campagnes avant correctifs} & \textbf{Scan \#5 (après correctifs)} \\
\midrule
Constat «~Panneau d'administration accessible sans authentification~» &
produit à chaque campagne (302 /wp-admin/ $\rightarrow$ wp-login.php suivie, 200 final) &
\textbf{0} \\
Sentinelle anti soft-404 & absente & exécutée (1 requête), tracée dans \code{recon\_json} \\
Avertissement «~cible bloque le scanner~» & absent (bilan présenté comme normal) &
affiché (racine 403, crawl quasi vide) \\
Potentiels faux positifs automatiques & au moins 1 (panneau admin) & \textbf{0} \\
\bottomrule
\end{tabular}
\caption{Effet des correctifs anti-faux-positifs mesuré sur la cible d'étude.}
\end{table}

\textbf{Analyse} : la cible restant bloquée (racine 403, crawl quasi vide --- cf. 10.4),
le scan \#5 ne produit que des constats passifs sur la racine (10 constats, re-vérifiés :
3 reproduits, 0 contredit). La mesure de corroboration reste celle du tableau 9.2 ;
l'apport démontré ici est la \textbf{disparition du faux positif récurrent} et la
transparence du bilan face au blocage. La mesure complète (avec crawl et corroboration
ZAP) sera refaite après levée du filtrage anti-bot.

% =====================================================================
\chapter{Difficultés rencontrées et solutions}

\section{Régression du matching de catégories (\code{rce} $\rightarrow$ «~resource~»)}

La comparaison en sous-chaîne faisait matcher le sigle \code{rce} dans
«~r\textbf{esou}rce~», recatégorisant les alertes «~Sub Resource Integrity~» en injection
de commande. \textbf{Solution} : matching par \textbf{mot entier} pour les sigles de
3 lettres, sous-chaîne tolérée à partir de 4 caractères (\code{cve-}, \code{.env}) ---
validé par les suites importer Nessus et ZAP.

\section{Décompression Brotli non supportée}

Décrit en 8.9.3 : cause racine d'un échec silencieux (corps illisible $\rightarrow$ 0
détection, crawl à 1 page). Leçon : \textbf{vérifier la chaîne complète} (transport
$\rightarrow$ décompression $\rightarrow$ parsing) avant d'incriminer la logique métier.

\section{Volume de sondes non borné sur les CMS}

Le crawl débouchant soudain sur 40 pages $\times$ 58 paramètres, le volume de sondes
explosait. \textbf{Solution} : bornes explicites (\code{MAX\_PROBE\_PAGES},
\code{MAX\_PARAMS\_PER\_SCAN}), sondes à coût réduit, arrêt anticipé --- la durée de scan
redevient prévisible (ES-02).

\section{Blocage anti-bot de la cible (limite de la mesure)}

Après plusieurs campagnes de scan, le WAF du site d'étude a fini par \textbf{dropper
silencieusement les connexions de l'adresse IP du scanner par fenêtres temporelles} ---
diagnostiqué par tests croisés (curl vs client interne, DNS, TCP, TLS). Les scans lancés
pendant une fenêtre de blocage échouent à la première requête racine (\code{ReconError}),
ce qui peut donner l'impression que «~plus rien n'est détecté~». \textbf{Solution} :
message d'erreur explicite anti-bot, attente entre campagnes, et documentation du
phénomène comme \textbf{limite méthodologique} : ZAP lui-même y est exposé.

\section{Sémantique des statuts (moteur vs analyste)}

La re-vérification posait initialement des statuts «~Confirmée~» automatiques --- en
contradiction avec la philosophie RF-12 (la confirmation est un \textbf{verdict humain}).
\textbf{Solution} : refonte RF-23 --- le moteur renforce la confiance (plafond 0,85),
l'analyste seul confirme (0,95) ; migration idempotente des anciennes données.

% =====================================================================
\chapter{Sécurité, éthique et conformité}

Le scan actif d'un tiers sans autorisation est illégal et éthiquement inacceptable.
\tscan{} intègre des garde-fous systémiques (checklist ES-01$\rightarrow$ES-12 du cahier
des charges) :

\begin{itemize}[leftmargin=1.5em]
  \item \textbf{ES-01} --- cible confirmée explicitement avant tout scan actif
        (\code{--authorized}) ;
  \item \textbf{ES-02} --- périmètre limitable : domaine, profondeur, durée maximale,
        types de tests ;
  \item \textbf{ES-03} --- mode sécurisé par défaut : sondes \textbf{GET bénignes
        uniquement} (charges à nonce, marqueurs dédiés) ; familles destructives exclues
        par construction (ES-04) ;
  \item \textbf{ES-05} --- journalisation systématique : chaque scan actif trace cible,
        horodatages, périmètre, chaque requête émise (\code{RequestHistory}), et les
        échecs de sondes ;
  \item \textbf{ES-06} --- traçabilité des corrections manuelles (\code{StatusHistory}
        avec raison) ;
  \item \textbf{ES-07} --- résultat brut d'origine conservé \textbf{sans modification} ;
  \item \textbf{ES-08/ES-09} --- preuves stockées localement ; aucune transmission à un
        tiers sans action explicite (lookups CVE/KEV à la demande, scan possible hors
        ligne) ;
  \item \textbf{ES-10} --- entrées validées avant construction des requêtes de test ;
  \item \textbf{ES-11} --- intégrité des mises à jour de connaissances vérifiée (taille
        max, format CVE) ;
  \item \textbf{ES-12} --- dépendances auditées (\code{pip-audit}).
\end{itemize}

Le protocole d'évaluation du chapitre 9 a respecté ces exigences : cible autorisée, volume
faible, mode non destructif, journalisation complète.

% =====================================================================
\chapter{Limites du MVP}

Les limites suivantes sont assumées et documentées pour l'encadrant (chapitre 15 du cahier
des charges) :

\begin{enumerate}[leftmargin=1.8em]
  \item \textbf{Familles non couvertes} : SSRF, désérialisation, IDOR approfondi ---
        exclues volontairement (ES-04, non destructivité) ;
  \item \textbf{Périmètre de crawl} : \tscan{} ne sonde pas les ressources binaires
        (images, PDF) --- certaines alertes ZAP sur ces ressources ne sont pas
        reproductibles par conception ;
  \item \textbf{Blocage anti-bot} : un WAF peut dégrader ou empêcher la mesure (fenêtres
        de drop) ; le moteur le signale mais ne le contourne pas ;
  \item \textbf{Le moteur ne confirme pas} : «~Confirmée~» reste un verdict d'analyste ---
        la corroboration apporte une indication, pas une preuve ;
  \item \textbf{Empaquetage} : l'exécutable Windows (PyInstaller/Inno Setup) a été
        produit le 28/09/2026 (deux exécutables one-file et un installeur
        \code{tscan-0.1.0-setup.exe}) ; reste la vérification d'installation sur machine
        «~propre~» ;
  \item \textbf{Tests manuels GUI UC1--UC5} : à exécuter sur poste avec affichage ;
  \item \textbf{Re-vérification limitée} : elle ne re-vérifie pas les alertes ZAP actives
        dont la logique de détection n'est pas implémentée côté \tscan{} (ex. injection
        dans un formulaire soumis) ;
  \item \textbf{Sentinelle et abstention} : le contrôle négatif anti soft-404 peut
        conduire à \textbf{s'abstenir} sur une exposition réelle si la ressource ressemble
        fortement à la page générique de la cible --- compromis assumé (RF-12) :
        l'incertitude produit une note pour revue analytique, jamais un constat ni une
        classification silencieuse ;
  \item \textbf{Corroboration «~historique~»} : la ré-observation d'un lot importé compte
        un scan \tscan{} antérieur sur la même cible comme preuve de recouvrement ---
        corroboration différée, pas un rejeu en temps réel de chaque alerte.
\end{enumerate}

% =====================================================================
\chapter{Perspectives post-stage}

\begin{itemize}[leftmargin=1.5em]
  \item \textbf{Installation sur machine «~propre~»} : tester l'installeur généré
        (exécutable Windows + installeur Inno Setup, produits le 28/09/2026) ;
  \item \textbf{Élargissement de la parité ZAP} : normalisation des scopes HTTP, alertes
        de session plus fines, alertes de type «~Content-Security-Policy~» sur les
        sous-ressources ;
  \item \textbf{Corroboration inversée} : utiliser la ré-observation pour
        \textbf{prioriser} les alertes importées non reproduites plutôt que pour les
        classer d'emblée en faux positifs ;
  \item \textbf{Base de connaissances} : enrichissement automatique des recommandations
        via CWE et le catalogue KEV (priorisation des vulnérabilités exploitées
        activement) ;
  \item \textbf{CI/CD} : pipeline GitHub Actions (pytest + ruff à chaque push) ;
  \item \textbf{Distribution} : publication des règles en paquet versionné indépendant du
        code ;
  \item \textbf{Contrôle négatif exhaustif} : la sentinelle anti soft-404, étendue en
        semaine 12 aux injections SQL, XSS, fichiers sensibles et listing de répertoire,
        doit encore couvrir les familles restantes (SSTI, injection de commande,
        traversée de chemins) dont les verdicts «~200~» restent exposés au même piège ;
  \item \textbf{Mesure précision / rappel} : corpus étiqueté (laboratoire + sites sains)
        pour chiffrer la précision du moteur avant/après correctifs et détecter
        objectivement les régressions ;
  \item \textbf{Statistiques de faux positifs par règle} : exploiter les verdicts
        d'analyste (\code{StatusHistory}, RF-12) pour produire un taux FP par règle et
        prioriser les améliorations --- sans auto-ajustement automatique des scores, qui
        ferait dériver le moteur.
\end{itemize}

% =====================================================================
\chapter{Bilan personnel et professionnel}

\section{Compétences développées}

\textbf{Techniques} : architecture logicielle (découplage cœur/CLI/GUI, modèle pivot),
sécurité web appliquée (OWASP Top 10 en pratique, charges bénignes, non-destructivité),
qualité logicielle (343 tests, linting, migrations), Python avancé (SQLAlchemy, httpx,
PySide6, Typer), diagnostic réseau/TLS (DNS, TCP, TLS, encodages de transport).

\textbf{Méthodologiques} : gestion de projet MVP, documentation vivante (checklist + README
reflétant l'état réel), traçabilité des décisions, communication régulière avec
l'encadrant, rédaction du cahier des charges.

\section{Apport pour un futur ingénieur en génie civil}

Bien que le stage relève de l'informatique, les compétences transférables au génie civil
sont réelles : \textbf{rigueur méthodologique} (protocoles de mesure, validation
expérimentale), \textbf{démarche qualité} (tests = contrôles de conformité), gestion de
projet structurée, et culture de la \textbf{sécurité des systèmes d'information},
désormais transverse à tout grand projet d'infrastructure (BIM, IoT sur chantiers, jumeaux
numériques, systèmes SCADA).

\section{Difficultés personnelles et apprentissages}

Le passage d'une logique «~ça fonctionne~» à une logique «~ça fonctionne \textbf{et} c'est
testé, mesuré, tracé~» a été l'apprentissage central. La gestion des échecs silencieux
(Brotli, blocage anti-bot) a appris à \textbf{déléguer le diagnostic à l'outil} (messages
explicites) plutôt qu'à l'utilisateur.

% =====================================================================
\chapter{Conclusion}

Le stage a permis de livrer \tscan{}, une plateforme fonctionnelle, testée et documentée
d'analyse, de corrélation et de validation de vulnérabilités web : import de 4 formats de
scanners, 54 règles YAML versionnées, 27 familles de détection non destructives, scoring
explicable, ré-observation avec corroboration mesurée (\textbf{76,9~\%} sur le cas
d'étude), reporting HTML/Markdown, CLI complète et application desktop --- le tout adossé
à une checklist de sécurité transverse (ES-01$\rightarrow$ES-12) et à \textbf{343 tests
automatisés}.

Au-delà de l'outil, le stage démontre qu'une \textbf{démarche d'ingénierie rigoureuse} ---
cahier des charges, MVP incrémental, validation expérimentale, limites assumées --- peut
produire, en 11 semaines, un système qui répond à une problématique réelle : la
fiabilisation des alertes de sécurité par corrélation et re-observation.

Les perspectives (installation sur machine «~propre~», CI/CD, élargissement de la parité
ZAP, corroboration inversée) tracent une feuille de route réaliste pour la suite du
produit après le stage.

% =====================================================================
\chapter{Références}

\begin{enumerate}[leftmargin=2em, label={[\arabic*]}]
  \item Cahier des charges \tscan{}, 16 chapitres --- \code{docs/cahier\_des\_charges.pdf} (projet).
  \item OWASP Foundation. \textit{OWASP Top 10:2021}. \url{https://owasp.org/Top10/}
  \item OWASP ZAP. \textit{Zed Attack Proxy --- User Guide}. \url{https://www.zaproxy.org/docs/}
  \item Nuclei (ProjectDiscovery). \textit{Template-based vulnerability scanner}.
        \url{https://docs.projectdiscovery.io/tools/nuclei/overview}
  \item Tenable. \textit{Nessus}. \url{https://www.tenable.com/products/nessus}
  \item Greenbone. \textit{OpenVAS}. \url{https://www.openvas.org/}
  \item NIST. \textit{National Vulnerability Database (NVD)}. \url{https://nvd.nist.gov/}
  \item CISA. \textit{Known Exploited Vulnerabilities (KEV) Catalog}.
        \url{https://www.cisa.gov/known-exploited-vulnerabilities-catalog}
  \item MITRE. \textit{Common Weakness Enumeration (CWE)}. \url{https://cwe.mitre.org/}
  \item NIST. \textit{Common Platform Enumeration (CPE)}. \url{https://cpe.mitre.org/}
  \item pytest. \textit{Framework de tests Python}. \url{https://docs.pytest.org/}
  \item httpx. \textit{HTTP client for Python}. \url{https://www.python-httpx.org/}
  \item SQLAlchemy. \url{https://www.sqlalchemy.org/}
  \item PySide6 / Qt for Python. \url{https://doc.qt.io/qtforpython/}
  \item Typer. \url{https://typer.tiangolo.com/}
  \item ruff. \textit{Linter Python}. \url{https://docs.astral.sh/ruff/}
  \item RFC 9110 --- \textit{HTTP Semantics}. \url{https://www.rfc-editor.org/rfc/rfc9110}
  \item RFC 7932 --- \textit{Brotli Compressed Data Format}. \url{https://www.rfc-editor.org/rfc/rfc7932}
\end{enumerate}

% =====================================================================
% ANNEXES
% =====================================================================
\appendix

\chapter{Commandes CLI principales}

\begin{table}[h!]
\centering
\small
\begin{tabular}{@{}lp{9.5cm}@{}}
\toprule
\textbf{Commande} & \textbf{Rôle} \\
\midrule
\code{tscan import <fichier>} & Importe un résultat Nuclei (JSONL), ZAP (JSON) ou Nessus/OpenVAS (XML) vers le modèle pivot \\
\code{tscan scan <cible>} & Lance un scan actif autorisé et non destructif sur la cible \\
\code{tscan correlate} & Corrèle et score les constats sur une même cible \\
\code{tscan list} & Liste les constats enregistrés avec filtres \\
\code{tscan correct} & Corrige manuellement le statut d'un constat (RF-12, ES-06) \\
\code{tscan report} & Génère un rapport HTML ou Markdown (RF-27--29) \\
\code{tscan lookup-cve <id>} & Interroge le cache local NVD pour une CVE \\
\code{tscan update-kev} & Synchronise le catalogue CISA KEV en cache local \\
\code{tscan rex <cible>} & Ré-observation : corrobore les imports avec un scan actif frais \\
\bottomrule
\end{tabular}
\caption{Commandes CLI principales de \tscan{}.}
\end{table}

\chapter{Rappel des exigences de sécurité (ES-01 à ES-12)}

\begin{itemize}[leftmargin=1.5em]
  \item \textbf{ES-01} : cible confirmée explicitement avant tout scan actif.
  \item \textbf{ES-02} : périmètre limitable (domaine, profondeur, durée, types de tests).
  \item \textbf{ES-03} : mode sécurisé (non destructif) activé par défaut.
  \item \textbf{ES-04} : familles à risque (SSRF, désérialisation\ldots) exclues du moteur.
  \item \textbf{ES-05} : journalisation systématique des scans actifs.
  \item \textbf{ES-06} : traçabilité des corrections manuelles de statut.
  \item \textbf{ES-07} : résultat brut d'origine conservé sans modification.
  \item \textbf{ES-08} : preuves stockées localement avec protection adaptée.
  \item \textbf{ES-09} : aucune transmission de données à un tiers sans action explicite.
  \item \textbf{ES-10} : entrées validées avant construction des requêtes de test.
  \item \textbf{ES-11} : intégrité vérifiée sur les données de mise à jour.
  \item \textbf{ES-12} : dépendances tierces suivies avec vigilance (\code{pip-audit}).
\end{itemize}

\end{document}
